\documentclass[../../../include/open-logic-section]{subfiles}

\begin{document}

\olfileid{sth}{choice}{tarskiscott}
\olsection{د تارسکي--سکاټ چاره}

په \olref[cardinals][cardsasords]{defcardinalasordinal} کښې مو
اصلي عددونه د ترتيبي عددونو په توګه تعريف کړي وو۔ د دې دپاره
مو د ښه ترتيب اصل فرض کړے و۔ دا کار مو يوازې ځکه وکړ چې دا
«عام منل شوے دود» دے۔

نو مخکې له دې چې د ښه ترتيب په اړه فلسفي مسئلې وڅېړو،
بايد روښانه کړو چې موږ له دغه عام دود څخه \emph{وتلے}
شو او د ښه ترتيب له فرضولو \emph{پرته} هم د اصلي عددونو
نظريه جوړولے شو۔ د $\cardeq{A}{B}$، $\cardle{A}{B}$
او $\cardless{A}{B}$ تعريفونه لا کارولے شو، لکه څنګه
چې په \olref[sfr][siz][]{chap} کښې راغلي دي۔ يوازې
د \emph{اصلي عدد} نوي تصور ته اړتيا لرو۔

يو ساده فکر دا دے چې د $A$ اصلي شمېر داسې تعريف کړو:
\[
	\Setabs{x}{\cardeq{A}{x}}.
\]
دا له هغې د $\# x Fx$ فرېګه‌يي تعريف سره پرتله کولے
شئ چې د \olref[cardinals][hp]{sec} په پاے کښې مو
په لنډه بڼه راوړے و۔ خو د هغو دلايلو له مخې چې هلته
مو ورته اشاره وکړه، دا تعريف نه سمېږي۔ هر يو غړيز
سټ له $\{\emptyset\}$ سره هم‌شمېره دے۔ خو د مراتبي
جوړښت په هر تالي پړاو کښې نوي يو غړيز سټونه جوړېږي
(يوازې د تېر پړاو يو غړيز سټ په پام کښې ونيسئ)۔
نو $\Setabs{x}{\cardeq{A}{x}}$ شتون نه لري، ځکه
د هغه لپاره يوه مرتبه نه شي ټاکل کېدے۔

د دې مسئلې د هوارولو دپاره د تارسکي او سکاټ
چاره کاروو:\footnote{يادونه: په ټولو فورمولونو کښې
پارامترونه کېدے شي، مګر که په ښکاره بل څه وويل شي۔}

\begin{defn}[تارسکي--سکاټ]
د هر فورمول $\phi(x)$ دپاره،
$[ x : \phi(x)] $ د ټولو هغو $x$ سټ وبولئ چې
تر ټولو ټيټه ممکنه مرتبه لري او $\phi(x)$ پوره کوي
(يا $\emptyset$، که هېڅ $\phi$ نۀ وي)۔
\end{defn}

بايد وګورو چې دا تعريف روا دے۔ په $\ZF$ کښې
\olref[spine][foundation]{zfentailsregularity} تضمينوي
چې $\setrank{x}$ د هر $x$ دپاره شتون لري۔ اوس که
څه داسې شيان وي چې $\phi$ پوره کوي، نو $\alpha$
هغه تر ټولو ټيټه مرتبه ټاکلے شو چې
$(\exists x\subseteq V_\alpha)\phi(x)$ وي؛ يعنې
$(\forall \beta
\in \alpha)(\forall x \subseteq V_\beta)\lnot \phi(x)$۔
بيا $[x : \phi(x)]$ د بېلونې په کارولو سره د
$\Setabs{x \in
V_{\alpha+1}}{\phi(x)}$ په توګه تعريفولے شو۔

د تارسکي--سکاټ چاره چې مو توجيه کړه، اوس يې د
اصلي شمېر د يوه تصور د تعريف دپاره کارولے شو:

\begin{defn}
د $A$ د \textsc{ts}-اصلي شمېر تعريف
$\text{tsc}(A) = [x :
\cardeq{A}{x}]$ دے۔
\end{defn}

د \textsc{ts}-اصلي عدد تعريف ښه ترتيب نه کاروي۔
خو بې له دې اصله هم ښودلے شو چې
\emph{\textsc{ts}-اصلي عددونه} د هغو
\emph{اصلي عددونو} په شان چلند کوي چې په
\olref[cardinals][cardsasords]{defcardinalasordinal}
کښې تعريف شوي دي۔ د بېلګې په توګه، که
\olref[cardinals][cardsasords]{lem:CardinalsBehaveRight}
او \olref[card-arithmetic][opps]{lem:SizePowerset2Exp}
د \textsc{ts}-اصلي عددونو په ژبه بيا بيان کړو،
ثبوتونه په $\ZF$ کښې هم، بې له ښه ترتيب فرضولو،
سم پر مخ ځي۔

په همدې ترڅ کښې، د تارسکي--سکاټ د چارې په وسيله
د ترتيبي عددونو نظريه هم جوړولے شو۔ که
$\tuple{A, <}$ ښه ترتيب وي، نو
$\text{tso}(A, <) = [\tuple{X,
R} : \ordeq{\tuple{A, <}}{\tuple{X, R}}]$
وبولئ۔ د اصلي او ترتيبي عددونو د دې لارې په اړه
نور په \citet[chs.~9--12]{Potter2004} کښې وګورئ۔

\end{document}